\section{A complex estimate with an additive length}\label{sec:complex}

The goal of this section is to bound the zeroth homology of a short
free complex from below using only the order of its differential
entries and the ranks of its terms. We first isolate the length
hypotheses, then prove the comparison between the complex and its
pullback to a blowup. A vector bundle with prescribed cohomological
roots will turn that comparison into the required inequality.

\subsection{Length hypotheses and the estimate}

All complexes in this section are indexed homologically. Thus a sheaf
cohomology group in degree $v$, arising from the term of a complex in degree
$i$, contributes to total degree $i-v$.

Let $C$ be a commutative ring, let $h\geq 1$, and put
$I=(z_1,\ldots,z_h)C$. Write $\mathscr T_I$ for the abelian category of
$C$-modules annihilated by some power of $I$. Suppose that
\[
 \lambda:\mathscr T_I\longrightarrow [0,\infty]
\]
is additive on short exact sequences and satisfies $\lambda(0)=0$.
We assume, in addition, that direct sums and filtered colimits of
length-zero modules annihilated by a common power of $I$ have length zero.
The standing numerical and homological assumptions are
\begin{equation}\label{eq:complex-length-hypotheses}
 \begin{split}
 0<\mu&:=\lambda(C/I)<\infty,\\
 \lambda\bigl(C/(z_1^a,\ldots,z_h^a)C\bigr)&=a^h\mu
       \qquad(a\geq 1),\\
 \lambda H_i(z_1^a,\ldots,z_h^a;C)&=0
       \qquad(a\geq 1,\ i>0).
 \end{split}
\end{equation}
No Noetherian hypothesis is imposed on $C$.

We say that a map between objects of $\mathscr T_I$ is an
\emph{isomorphism modulo length zero} if its kernel and cokernel have
length zero. The length-zero objects form a Serre subcategory. Consequently
one may work in the quotient by this subcategory when taking kernels,
images, cokernels, or finite spectral sequences. In particular, a
commutative square whose two vertical maps are isomorphisms modulo length
zero induces such an isomorphism between the images of its horizontal
maps. All uses of this convention below involve finite filtrations.

We will also use the following elementary observation about complexes
with finite exhaustive filtrations in a common index range. If a
filtration-preserving map induces zero on the first page of its spectral
sequence, it induces zero on the associated graded of homology and hence
lowers the filtration on homology by one. A composition of as many such
maps as there are filtration steps is therefore zero on homology. The same
statement holds modulo length zero whenever the pages under consideration
belong to $\mathscr T_I$.

\begin{theorem}\label{thm:complex}
Let $C$, $I$, and $\lambda$ be as above. Assume that $\lambda$ is additive
on short exact sequences, that $\lambda(0)=0$, and that direct sums and
filtered colimits of length-zero modules killed by a common power of $I$
have length zero. Assume also \eqref{eq:complex-length-hypotheses}.
Let $F$ be a complex of finite
free $C$-modules, concentrated in degrees $0,\ldots,h$, which is
contractible locally on $\Spec C\setminus V(I)$. Set
$b_i=\rank_C F_i$. Then $H_i(F)$ is annihilated by a power of $I$,
\[
 \lambda H_i(F)=0\quad(i>0),\qquad \lambda H_0(F)<\infty.
\]
Suppose further that
\[
 \sum_{i=0}^h(-1)^i b_i=0
 \quad\text{and that all differential entries belong to }I^s,
 \qquad s\geq1.
\]
There is a constant $A_h$, depending only on $h$, such that
\begin{equation}\label{eq:complex-estimate}
 \lambda H_0(F)\ \geq\
 \mu s^h\left(b_0-\frac{A_h}{s}\sum_{i=0}^h b_i\right).
\end{equation}
In particular, for fixed $h$ and fixed ranks the error is $O_h(1/s)$,
uniformly in $C$, $I$, $\lambda$, and $F$.
\end{theorem}

\subsection{Koszul acyclicity modulo length zero}

We first prove the homology assertion, which does not require the
conditions on ranks or differential entries.

\begin{lemma}\label{lem:acyclicity}
Under the standing length hypotheses, let $F$ be a finite free
complex in degrees $0,\ldots,h$, contractible locally off $V(I)$.
Each $H_i(F)$ is annihilated by
a fixed power of $I$, all its lengths are finite, and its positive
homology has length zero.
\end{lemma}
\begin{proof}
For each $j$, the complex $F[1/z_j]$ is exact. A bounded exact complex of
projective modules is contractible, so clearing the finitely many
denominators in a contracting homotopy shows that some power of $z_j$
acts nullhomotopically on $F$. Choose a common exponent $a_0$.
It follows that $(z_1^{a_0},\ldots,z_h^{a_0})$ annihilates every $H_i(F)$,
and therefore so does $I^{h(a_0-1)+1}$.

Put $K_a=K(z_1^a,\ldots,z_h^a;C)$. If $a\geq a_0$, each multiplication
by $z_j^a$ is nullhomotopic on $F$ and also on its tensor product with any
of the other Koszul factors. Splitting the resulting mapping cones gives
\begin{equation}\label{eq:complex-koszul-splitting}
 F\otimes_C K_a\simeq
 \bigoplus_{T\subseteq\{1,\ldots,h\}}F[|T|].
\end{equation}
The splittings can be chosen to preserve the inclusion of the copy of
$F$ in Koszul degree zero.

The homology of $F\otimes K_a$ is annihilated by
$(z_1^a,\ldots,z_h^a)$, since these multiplications are nullhomotopic on
$K_a$. In the bounded double-complex spectral sequence obtained by taking
Koszul homology first, all positive Koszul rows have length zero by
\eqref{eq:complex-length-hypotheses}. The remaining row has terms
\[
 F_i\otimes_C C/(z_1^a,\ldots,z_h^a),\qquad 0\leq i\leq h,
\]
of finite length. Hence every $H_n(F\otimes K_a)$ has finite length, and
its length is zero when $n>h$. The summand $F[h]$ in
\eqref{eq:complex-koszul-splitting} proves the assertions: $H_0(F)$ is a
summand of $H_h(F\otimes K_a)$, and $H_i(F)$ for $i>0$ is a summand of
$H_{h+i}(F\otimes K_a)$.
\end{proof}

\subsection{The blowup and its exceptional divisor}

The homology lengths are now finite. To estimate them, we pass to a
space on which $I$ becomes invertible, while keeping track of the
Euler characteristic of the original complex. Let
\[
 Y=\Proj_C\Bigl(\bigoplus_{n\geq0}I^n\Bigr),\qquad
 L=I\cO_Y=\cO_Y(1),\qquad E=V(L).
\]
Here $\cO_Y(1)$ is the Rees-algebra twist. Indeed, the standard chart
corresponding to $z_i$ has ring $C[I/z_i]\subseteq C[1/z_i]$. On this
chart $I$ becomes $(z_i)$ and $z_i$ is a nonzerodivisor. Thus $L$ is an
invertible ideal, and $E$ is an effective Cartier divisor, even when $C$
is not Noetherian. Empty charts may simply be omitted. The generators of
$I$ give a closed immersion $Y\longrightarrow\PP_C^{h-1}$, and
\[
 E=\Proj(\gr_I C)\longrightarrow\PP_{C/I}^{h-1}.
\]
For a vector bundle $V$ on $\PP_\ZZ^{h-1}$ we use the same symbol for its
pullbacks to these schemes, and write $V(t)=V\otimes\cO(t)$.

\begin{lemma}\label{lem:complex-graded-comparison}
The graded surjection
\begin{equation}\label{eq:complex-graded-map}
 (C/I)[Z_1,\ldots,Z_h]\longrightarrow\gr_I C
\end{equation}
has length-zero kernel in every degree. For every vector bundle $V$ on
$\PP_\ZZ^{h-1}$ and every integer $t$, restriction induces isomorphisms
modulo length zero
\[
 H^v(\PP_{C/I}^{h-1},V(t))\longrightarrow H^v(E,V(t)).
\]
All of these cohomology modules have finite length. Moreover, for $t\geq0$
the natural map
\[
 \gr_I^t C\longrightarrow H^0(E,\cO_E(t))
\]
is an isomorphism modulo length zero, and $H^v(E,\cO_E(t))$ has length
zero for $v>0$.
\end{lemma}
\begin{proof}
Set $J_a=(z_1^a,\ldots,z_h^a)$. Since
$I^{h(a-1)+1}\subseteq J_a$, the $I$-adic filtration of $C/J_a$ is finite.
There is a graded surjection
\[
 (C/I)[Z_1,\ldots,Z_h]/(Z_1^a,\ldots,Z_h^a)
 \longrightarrow\gr_I(C/J_a).
\]
The sum of the lengths of the homogeneous pieces on each side is $a^h\mu$.
Additivity and nonnegativity imply that every homogeneous piece of its
kernel has length zero. For a given degree $n$, choose $a>n$. The relations
$Z_i^a$ have no effect in degree $n$, and $J_a\subseteq I^{n+1}$, so the
degree-$n$ map is precisely \eqref{eq:complex-graded-map}. This proves the
first assertion. It also gives
\begin{equation}\label{eq:complex-power-colength}
 \lambda(C/I^M)=\mu\binom{M+h-1}{h}\qquad(M\geq1).
\end{equation}

Write $P=(C/I)[Z_1,\ldots,Z_h]$ and let $N$ be the kernel in
\eqref{eq:complex-graded-map}. The degree-zero component of a homogeneous
localization of $N$, as used on any standard projective chart or its
intersections, is a filtered colimit of homogeneous pieces of $N$.
All these pieces are annihilated by $I$ and have length zero. The assumed
colimit property therefore gives length zero for every such localized
component. Tensoring with the restriction of $V(t)$ preserves this
property, since that restriction is finite projective over the chart ring.
The finite standard affine Cech complexes for projective space and $E$
consequently have termwise length-zero kernels, and restriction is an
isomorphism modulo length zero on their cohomology.

To check finiteness without a hypothesis on $C/I$, resolve $V(t)$ on
$\PP_\ZZ^{h-1}$ by finitely many finite sums of twists of the structure
sheaf. Such a resolution exists by taking a finite graded free resolution
of a finitely generated graded module representing $V(t)$ over
$\ZZ[Z_1,\ldots,Z_h]$. This polynomial ring is regular of finite
dimension and therefore has finite global dimension
\cite[Tags~07NF and~00OE]{Stacks};
its finitely generated graded projectives are graded free, as follows by
lifting a homogeneous basis modulo $(Z_1,\ldots,Z_h)$ and applying graded
Nakayama. The sheaf resolution is locally split because its final cokernel
is a vector bundle, so it remains exact after arbitrary base change.
The cohomology of every twist on $\PP_{C/I}^{h-1}$ is finite free over
$C/I$. A finite spectral sequence for the resolution proves the required
finiteness. Finally, for $t\geq0$ the polynomial degree-$t$ piece is the
space of global sections of $\cO(t)$ and its higher cohomology vanishes.
Combining these facts with \eqref{eq:complex-graded-map} proves the last
assertions.
\end{proof}

\subsection{Euler characteristics on the blowup}

We now compare the homology lengths of $F$ with those of its pullback
to $Y$. We will also show that, when the alternating rank is zero,
tensoring the pullback with a vector bundle multiplies its Euler
characteristic by the bundle's rank.

For a bounded complex $D$ of quasi-coherent sheaves on $Y$, write
\[
 \mathbb H_k(Y,D)=H_k\mathbf R\Gamma(Y,D).
\]
These groups can be calculated using the finite standard affine Cech
cover, whose intersections are affine. Whenever they are $I$-power torsion
of finite length, put
\[
 \chi_\lambda(D)=\sum_k(-1)^k\lambda\mathbb H_k(Y,D).
\]
Only finitely many degrees occur in our applications.

Pull $F$ back to $Y$ and denote the resulting complex by $F_Y$. If the
differential entries belong to $I^s$, define a complex $G$ by
\begin{equation}\label{eq:complex-G}
 G_i=L^{-is}\otimes_C F_i=\cO_Y(-is)^{b_i}.
\end{equation}
Its differential is induced by that of $F_Y$ after inverting $L$.
On a chart where $L=(z_j)$, the matrix in the bases
$z_j^{-is}F_i$ is $z_j^{-s}d_i$, where $d_i$ is the matrix for $F$.
It is regular because the entries of $d_i$ belong to $L^s$.
The natural inclusions $F_i\to L^{-is}\otimes_C F_i$ give an
inclusion of complexes $F_Y\subseteq G$.

In degree $i$, the quotient $G/F_Y$ has successive layers
$\cO_E(l)^{b_i}$ for $l=-is,\ldots,-1$. When the alternating rank is
zero, these finite layer ranges will account, after the Euler comparison,
for the entire difference between $\chi_\lambda(G\otimes V)$ and
$(\rank V)\lambda H_0(F)$. We first justify the comparison for an
arbitrary bundle $V$; the choice of $V$ that gives a nonnegative Euler
characteristic will follow.

\begin{lemma}\label{lem:complex-twist-maps}
Let $D$ be either $F_Y\otimes V$ or $G\otimes V$, with $V$ a vector
bundle pulled back from $\PP_\ZZ^{h-1}$. There is an integer $M_D$ such
that, for every integer $j$ and every $M\geq M_D$, the inclusion
\[
 D(j+M)\longrightarrow D(j)
\]
induces zero on hyperhomology. All hyperhomology modules of $D(j)$ are
annihilated by a power of $I$ and have finite length.
\end{lemma}
\begin{proof}
On a standard affine chart $U_i$, $L$ is generated by the nonzerodivisor
$z_i$. The chart complex of $D$ consists of finite projective modules and
becomes contractible after inverting $z_i$: away from $E$, both $F_Y$ and
$G$ identify with the pullback of $F$. Clearing denominators in a
contracting homotopy gives an exponent $m_i$ for which multiplication by
$z_i^{m_i}$ is nullhomotopic on this chart. Trivializing $L$ shows that the
same exponent works for every twist. It also works on all intersections
contained in this chart. Choose $m$ at least as large as every $m_i$.

The map $D(j+m)\to D(j)$ is zero on the internal-homology page of the
Cech double complex, since on each intersection it is identified with a
nullhomotopic multiplication by a local equation to the power $m$.
The Cech filtration has at most $h$ steps. The filtered-complex observation
above therefore shows that $D(j+hm)\to D(j)$ is zero on hyperhomology.
Every larger twist map factors through one of these, so $M_D=hm$ works.

For $a\in I^{M_D}$, multiplication by $a$ on $D(j)$ factors as
\[
 D(j)\xrightarrow{a}D(j+M_D)\longrightarrow D(j).
\]
Thus $I^{M_D}$ annihilates its hyperhomology. If $M\geq M_D$, the long
exact sequence for
\[
 0\longrightarrow D(j+M)\longrightarrow D(j)
 \longrightarrow D(j)/D(j+M)\longrightarrow0
\]
injects $\mathbb H_k(Y,D(j))$ into the hyperhomology of the quotient.
That quotient has a finite filtration by $E$-layers, each of whose terms
is a finite sum of bundles $V(t)|_E$. Their sheaf cohomology has finite
length by Lemma~\ref{lem:complex-graded-comparison}. The two finite
spectral sequences, first for each layer and then for the filtration,
give finite length for the quotient's hyperhomology and hence for that
of $D(j)$.
\end{proof}

The Euler characteristics of finite thickenings need not recover that of $F$.
For example, in the alternating-rank-zero case needed for the estimate,
\[
 \sum_k(-1)^k\lambda H_k(F/I^M F)
 =\lambda(C/I^M)\sum_i(-1)^i b_i=0.
\]
Instead, we compare the images of homology maps from a sufficiently
large thickening to a fixed smaller one. These images will recover the
homology of $F$ and of $F_Y$, modulo length zero. This argument gives
the first assertion below without the alternating-rank hypothesis;
that hypothesis is used only afterward to prove invariance under twists
and the vector-bundle formula.

\begin{lemma}\label{lem:complex-euler-comparison}
One has
\begin{equation}\label{eq:complex-euler-comparison}
 \chi_\lambda(F_Y)=\sum_i(-1)^i\lambda H_i(F)=\lambda H_0(F).
\end{equation}
If $\sum_i(-1)^i b_i=0$, then for every vector bundle $V$ pulled back from
$\PP_\ZZ^{h-1}$,
\begin{equation}\label{eq:complex-bundle-rank}
 \chi_\lambda(F_Y\otimes V)=(\rank V)\chi_\lambda(F_Y).
\end{equation}
\end{lemma}
\begin{proof}
\emph{Comparison on finite thickenings.}
For each $M\geq1$ there is a natural comparison
\begin{equation}\label{eq:complex-thick-comparison}
 F/I^M F\longrightarrow
 \mathbf R\Gamma(Y,F_Y/L^M F_Y).
\end{equation}
Filter its two sides by powers of $I$ and $L$, respectively. At the
$l$th layer, $0\leq l<M$, comparison for each free summand is
\[
 \gr_I^l C\longrightarrow\mathbf R\Gamma(E,\cO_E(l)).
\]
Lemma~\ref{lem:complex-graded-comparison} says that this is an isomorphism
on homology modulo length zero. The finite filtration, together with the
finite complex $F$, proves the same assertion for
\eqref{eq:complex-thick-comparison}. These comparisons commute with all
transition maps as $M$ increases.

\smallskip
\noindent\emph{The geometric tower.}
We identify the eventual images of these two towers. On the geometric
side put
\[
 A=\mathbf R\Gamma(Y,F_Y),\qquad
 B_M=\mathbf R\Gamma(Y,F_Y/L^M F_Y).
\]
The triangles
\begin{equation}\label{eq:complex-thick-triangle}
 A\longrightarrow B_M\longrightarrow
 \mathbf R\Gamma(Y,F_Y\otimes L^M)[1]\longrightarrow A[1]
\end{equation}
are compatible with the transition $N\to M$ for $N\geq M$.
Lemma~\ref{lem:complex-twist-maps} implies that $H_k(A)\to H_k(B_M)$ is
injective once $M$ is sufficiently large. It also says that the transition
on the third terms of \eqref{eq:complex-thick-triangle} is zero on homology
once $N-M$ is sufficiently large. Exactness then gives
\begin{equation}\label{eq:complex-geometric-image}
 \operatorname{im}\bigl(H_k(B_N)\to H_k(B_M)\bigr)
   =\operatorname{im}\bigl(H_k(A)\to H_k(B_M)\bigr)
   \cong H_k(A).
\end{equation}
The reverse inclusion in this equality uses only the compatible maps from
$A$ to all $B_N$.

\smallskip
\noindent\emph{The Koszul tower.}
For the algebraic side first use $J_a=(z_1^a,\ldots,z_h^a)$ and
$A_a=F\otimes K_a$. The Koszul augmentation gives a compatible map
\[
 A_a\longrightarrow F/J_aF
\]
which is an isomorphism on homology modulo length zero. Indeed its cone
is obtained by tensoring the bounded free complex $F$ with the cone of
$K_a\to C/J_a$; all homology of the latter cone has length zero by
\eqref{eq:complex-length-hypotheses}.

The inclusion $F\subseteq A_a$ is the subcomplex in Koszul degree zero.
For $a\geq a_0$ it is split on homology by
\eqref{eq:complex-koszul-splitting}. Let $Q_a=A_a/F$. The standard
transition $K_b\to K_a$, for $b\geq a$, is the identity in degree zero
and multiplies the basis vector indexed by $T$ by
$\prod_{j\in T}z_j^{b-a}$. Filter $Q_a$ increasingly by Koszul
degrees $1,\ldots,p$, for $p=1,\ldots,h$.
Its associated graded complexes are copies of shifts of $F$, and the
transition on each such piece has a nonempty index set $T$.
Choose $t$ with $I^tH(F)=0$, as in Lemma~\ref{lem:acyclicity}.
For $b-a\geq t$ every transition is zero on the homology of these graded
pieces. A composition of $h$ such transitions is therefore zero on
$H(Q_a)$. For example, $H(Q_b)\to H(Q_a)$ is zero when $b-a\geq ht$.
The long exact sequences for $0\to F\to A_a\to Q_a\to0$ now give,
for $a$ sufficiently large and then $b$ sufficiently large,
\[
 \operatorname{im}\bigl(H_k(A_b)\to H_k(A_a)\bigr)
  =\operatorname{im}\bigl(H_k(F)\to H_k(A_a)\bigr)
  \cong H_k(F).
\]
Passing through the Koszul augmentations proves this assertion modulo
length zero for the tower $F/J_aF$.

\smallskip
\noindent\emph{From Koszul powers to ordinary powers.}
Here is explicitly why it passes to the ordinary-power tower
$P_M=F/I^M F$. The containments
\begin{equation}\label{eq:complex-cofinality}
 I^{h(a-1)+1}\subseteq J_a\subseteq I^a
\end{equation}
give all the required quotient maps. Choose one sufficiently large $a$.
For $M\geq h(a-1)+1$, the composite
$H_k(F)\to H_k(P_M)\to H_k(F/J_aF)$ proves injectivity of the first
map modulo length zero. Next fix such an $M$, choose $a\geq M$ large,
and choose $b\geq a$ so that the image from level $b$ to level $a$ in
the $J$-tower is the image of $H_k(F)$. For $N\geq h(b-1)+1$, the
transition factors through
\[
 P_N\longrightarrow F/J_bF\longrightarrow F/J_aF\longrightarrow P_M.
\]
Its homology image is consequently contained in the image of $H_k(F)$
modulo length zero; compatibility gives the opposite containment. Thus
the eventual image in the $P$-tower identifies with $H_k(F)$ modulo
length zero.

Choose $M$ and then $N$ large enough for both this description and
\eqref{eq:complex-geometric-image}. The compatible comparisons
$P_N\to B_N$ and $P_M\to B_M$ in
\eqref{eq:complex-thick-comparison} identify the homology images of
$P_N\to P_M$ and $B_N\to B_M$ modulo length zero. These images
identify with $H_k(F)$ and $H_k(A)$, respectively, modulo length zero.
All their lengths are finite by
Lemma~\ref{lem:acyclicity}, \eqref{eq:complex-power-colength}, and
Lemma~\ref{lem:complex-twist-maps}. Therefore
\[
 \lambda\mathbb H_k(Y,F_Y)=\lambda H_k(F)
\]
for every $k$, which proves \eqref{eq:complex-euler-comparison}.

\smallskip
\noindent\emph{Tensoring by a vector bundle.}
The homology comparison is established. To prove
\eqref{eq:complex-bundle-rank}, first compare consecutive twists
of $F_Y$. The quotient of $F_Y(j)$ by $F_Y(j+1)$ has term
$\cO_E(j)^{b_i}$ in degree $i$, so its Euler characteristic is
\[
 \left(\sum_i(-1)^i b_i\right)
 \sum_v(-1)^v\lambda H^v(E,\cO_E(j))=0.
\]
Euler additivity shows that $\chi_\lambda(F_Y(j))$ is independent of
$j\in\ZZ$. Finally resolve $V$ on $\PP_\ZZ^{h-1}$ by a finite complex
of finite sums of twists, as in
Lemma~\ref{lem:complex-graded-comparison}. The resolution remains exact
on $Y$. Tensoring with $F_Y$ and taking Euler characteristics expresses
$\chi_\lambda(F_Y\otimes V)$ as the alternating sum of the corresponding
twist characteristics. The alternating sum of the ranks of the resolution
is $\rank V$, giving \eqref{eq:complex-bundle-rank}.
\end{proof}

\subsection{Bundles with prescribed cohomological roots}

We next construct the bundle that determines the signs of the needed
Euler characteristic. We use the product-of-projective-lines
construction of Eisenbud--Schreyer
\cite[Theorem~6.1]{EisenbudSchreyer}, verifying its integral and
base-change properties so that it applies over our auxiliary ring.

\begin{lemma}\label{lem:complex-root-bundle}
Given integers $r_1>\cdots>r_{h-1}$, there is a vector bundle $V$ of rank
$(h-1)!$ on $\PP_\ZZ^{h-1}$ with the following property. For every ring
$A$ and every integer $a$, $V(a)$ on $\PP_A^{h-1}$ is acyclic if $a=r_j$
for some $j$. Otherwise its cohomology is a finite free $A$-module in
the single degree
\[
 v(a)=\#\{j:r_j>a\},
\]
and its Euler characteristic in ranks is
\[
 \prod_{j=1}^{h-1}(a-r_j).
\]
For $A=C/I$, the Euler characteristic with respect to $\lambda$ is
$\mu\prod_j(a-r_j)$, and the same Euler characteristic and cohomology
vanishing statements hold on $E$, modulo length zero.
\end{lemma}
\begin{proof}
Put $n=h-1$. If $n=0$, use the rank-one bundle on $\PP_\ZZ^0$.
For $n>0$, multiplication of binary linear forms defines
\[
 f:(\PP_\ZZ^1)^n\longrightarrow\PP_\ZZ^n,
 \qquad f^*\cO(1)=\cO(1,\ldots,1).
\]
The map is proper and has finite geometric fibers, since a binary form
has only finitely many ordered factorizations into linear factors, up to
the scalars already removed by the projective coordinates. Hence $f$ is
finite. On every field fiber, source and target are smooth of dimension
$n$. Corresponding local rings have equal dimensions, since the residue
field extension is finite and the dimensions are $n$ minus the respective
residue-field transcendence degrees. Thus the source is Cohen--Macaulay
and satisfies the dimension criterion for flatness over the regular
target~\cite[Tag~00R4]{Stacks}. This gives flatness on each field fiber.
Since both schemes are
flat over $\ZZ$, the
fiberwise flatness criterion~\cite[Tag~039D]{Stacks} gives flatness
of $f$. Over $\QQ$ the finite
extension of function fields is separable. A geometric generic form has
$n$ distinct factors and exactly $n!$ ordered factorizations. Its reduced
geometric generic fiber therefore has $n!$ points, so the generic degree,
and hence the constant rank of $f$, is $n!$.

Define
\[
 V=f_*\cO(-1-r_1,\ldots,-1-r_n).
\]
Finite flatness makes this a vector bundle of rank $n!$, and its formation
commutes with arbitrary base change. The projection formula reduces the
cohomology of $V(a)$ to that of
$\cO(a-1-r_1,\ldots,a-1-r_n)$ on $(\PP_A^1)^n$.
On $\PP_A^1$, the bundle $\cO(m)$ has free cohomology only in degree zero
for $m\geq0$, only in degree one for $m\leq-2$, and none for $m=-1$;
its Euler characteristic is $m+1$. These computations commute with any
base change and their complexes split over $\ZZ$. Tensoring the complexes
for the $n$ factors proves the asserted cohomological degree and product
formula over every $A$. For $A=C/I$, a free module of rank $q$ has length
$q\mu$. The statements on $E$ follow from
Lemma~\ref{lem:complex-graded-comparison}.
\end{proof}

\subsection{The sign and the leading coefficient}

\begin{proof}[Proof of Theorem~\ref{thm:complex}]
The homology assertions are Lemma~\ref{lem:acyclicity}. It remains to
prove the estimate. We choose the bundle roots to serve two purposes:
their positions force the hyperhomology of $G\otimes V$ into degree zero,
modulo length zero, while their limiting values make every layer sum
with $1\leq i\leq h$ have the same leading term $(h-1)!s^h$.
The alternating-rank identity will then isolate $b_0$.
The derivative of the following polynomial provides exactly these roots:
\[
 H(x)=\prod_{i=1}^h(x+i),
\]
and write $\rho_1>\cdots>\rho_{h-1}$ for the roots of $H'(x)$.
Rolle's Theorem gives
\[
 -(j+1)<\rho_j<-j.
\]
For all sufficiently large $s$, depending only on $h$, round $s\rho_j$
to integers $r_j$ so that
\begin{equation}\label{eq:complex-root-positions}
 -(j+1)s<r_j<-js,\qquad |r_j-s\rho_j|\leq1.
\end{equation}
Let $V$ be the bundle of Lemma~\ref{lem:complex-root-bundle}, and put
$D=G\otimes V$ with $G$ as in \eqref{eq:complex-G}.

We claim that
\begin{equation}\label{eq:complex-positive-chi}
 \chi_\lambda(D)\geq0.
\end{equation}
In the $l$th $E$-layer of a twist of $D$, the term in homological degree
$i$ is $V(l-is)|_E^{b_i}$. If $l\geq0$, then
\eqref{eq:complex-root-positions} implies
\[
 \#\{j:r_j>l-is\}\leq i.
\]
Indeed a root with index $j\geq i$ is strictly less than $-is$, and
for $i=0$ all roots are negative. Thus the layer's sheaf cohomology in
degree $v>i$ has length zero. In the finite hyperhomology spectral
sequence no term of negative total degree $i-v$ survives modulo length
zero. Filtering by the layers $l=0,\ldots,M-1$ shows that
\[
 \lambda\mathbb H_k(Y,D/D(M))=0\qquad(k<0).
\]
For $M$ large, Lemma~\ref{lem:complex-twist-maps} makes the map on
hyperhomology from $D(M)$ to $D$ zero. The quotient long exact sequence
therefore injects $\mathbb H_k(Y,D)$ into
$\mathbb H_k(Y,D/D(M))$, proving length zero for $k<0$.

If $l<0$, then
\[
 \#\{j:r_j>l-is\}\geq i-1.
\]
For $i\geq1$ this follows because every $j\leq i-1$ has
$r_j>-(j+1)s\geq-is>l-is$; for $i=0$ the assertion is vacuous.
Consequently the corresponding layer has hyperhomology of length zero
in total degrees greater than $1$. Filtering the quotient with layers
$l=-M,\ldots,-1$ gives
\[
 \lambda\mathbb H_k(Y,D(-M)/D)=0\qquad(k>1).
\]
For large $M$ the map from $D$ to $D(-M)$ is zero on hyperhomology.
The long exact sequence gives a surjection
\[
 \mathbb H_{k+1}(Y,D(-M)/D)\longrightarrow\mathbb H_k(Y,D),
\]
and hence length zero for $k>0$. Lemma~\ref{lem:complex-twist-maps}
supplies finite length in all degrees. We have proved
$\chi_\lambda(D)=\lambda\mathbb H_0(Y,D)\geq0$, as claimed.

The inclusion $F_Y\otimes V\subseteq G\otimes V$ has, in degree $i$,
quotient $(L^{-is}/\cO_Y)\otimes V^{b_i}$, with successive layers
$V(l)|_E^{b_i}$ for $l=-is,\ldots,-1$. All Euler characteristics here
are finite. Additivity, Lemma~\ref{lem:complex-euler-comparison}, and
Lemma~\ref{lem:complex-root-bundle} give
\begin{equation}\label{eq:complex-sum-estimate}
 0\leq (h-1)!\lambda H_0(F)
 +\mu\sum_{i=0}^h(-1)^i b_i
       \sum_{l=-is}^{-1}\prod_{j=1}^{h-1}(l-r_j).
\end{equation}
The sum for $i=0$ is empty.

We make the error in these sums explicit. Set
\[
 P_s(x)=\prod_{j=1}^{h-1}(x-r_j/s),\qquad
 P(x)=\prod_{j=1}^{h-1}(x-\rho_j)=\frac{H'(x)}h.
\]
On the fixed interval $[-h,0]$, \eqref{eq:complex-root-positions}
implies $\|P_s-P\|_\infty=O_h(s^{-1})$, while $P$ and its derivative
are bounded in terms of $h$ alone. The elementary error estimate for
left-endpoint Riemann sums consequently gives, uniformly for
$i=1,\ldots,h$,
\begin{align*}
 s^{-h}\sum_{l=-is}^{-1}\prod_j(l-r_j)
 &=\frac1s\sum_{l=-is}^{-1}P_s(l/s)\\
 &=\int_{-i}^0P(x)\,dx+O_h(s^{-1})\\
 &=\frac{H(0)-H(-i)}h+O_h(s^{-1})
  =(h-1)!+O_h(s^{-1}).
\end{align*}
Since $\sum_{i=1}^h(-1)^i b_i=-b_0$, substitution into
\eqref{eq:complex-sum-estimate} yields
\[
 \lambda H_0(F)\geq
 \mu s^h\left(b_0-O_h\left(\frac{\sum_i b_i}{s}\right)\right).
\]
The error bound and the threshold for $s$ depend only on $h$.
The homotopy exponents and truncation levels used earlier may depend on
$C$, $F$, and $V$; they establish exact comparisons and vanishing modulo
length zero and do not enter this numerical error. This proves
\eqref{eq:complex-estimate} for all sufficiently large $s$. Increase
$A_h$, if necessary, to be at least the threshold for $s$. For the
remaining positive integers $s$, its right side is nonpositive because
$b_0\leq\sum_i b_i$, so nonnegativity of length proves the same inequality.
For $h=1$ the empty-product conventions give the argument directly,
with exact sums and no error. The proof is complete.
\end{proof}
